\subsection{Зачем нужен EDA для моделирования, примеры}

\textbf{Цель EDA} -- понять структуру данных, выявить закономерности, аномалии и взаимосвязи \textbf{до}
построения моделей.
Это критический этап, который напрямую влияет на качество и интерпретируемость итоговой модели.

\begin{itemize}
    
    \item \textbf{Понимание данных и генерация гипотез:} EDA помогает сформулировать осмысленные гипотезы о
    взаимосвязях между переменными.
    Без этого моделирование может стать «слепым» поиском.
    
    \item \textbf{Примеры обнаружения и обработки проблем с данными}
    
    \begin{itemize}
        
        \item \textit{Пропущенные значения:} при анализе медицинских данных EDA может показать, что в 40 \% записей
        отсутствует показатель уровня сахара.
        Это требует решения: удалить строки, импутировать значения или добавить бинарный признак «пропуск».
        
        \item \textit{Выбросы (аномалии):} в данных о транзакциях EDA (например, через box-plot) может выявить
        платежи на сумму \$10 млн при медиане в \$50.
        Это может быть ошибка ввода, мошенничество или реальная транзакция, требующая отдельного анализа.
    
    \end{itemize}
    
    \item \textbf{Проверка предположений модели:} многие алгоритмы (линейная регрессия, логистическая регрессия)
    строятся на статистических предположениях
    Например, построив гистограмму и Q-Q plot целевой переменной «цена дома», можно обнаружить её сильную
    правостороннюю асимметрию.
    Применение логарифмического преобразования (\textit{log-transform}) сделает распределение ближе к
    нормальному и повысит эффективность линейной модели.
    
    \item \textbf{Выбор признаков и их инжиниринг:} EDA выявляет избыточные или слабые признаки.
    Наприме, матрица корреляций может показать почти линейную зависимость (\textit{multicollinearity}) между
    признаками «площадь дома в квадратных метрах» и «площадь в квадратных футах».
    Один из них следует исключить.
    
    \item \textbf{Визуальное представление для заказчиков:} графики EDA (гистограммы, scatter plots) — лучший
    способ донести факты о данных до заказчика или команды, не погружённой в технические детали

\end{itemize}

\subsection{Типичные шаги EDA}

Процесс EDA является итеративным и не строго линейным.

\begin{enumerate}
    
    \item \textbf{Первичный осмотр данных}
    
    \begin{itemize}
        \item Загрузка данных (\texttt{pandas.read\_csv()} и т.д.)
        \item \texttt{.head()}, \texttt{.tail()} — просмотр первых/последних строк
        \item \texttt{.info()} -- информация о типах данных и наличии пропусков.
        \item \texttt{.describe()} -- сводная статистика (среднее, std, минимум и далее) для числовых признаков.
        \item \texttt{.shape} — размерность данных (число строк и столбцов)
    \end{itemize}
    
    \item \textbf{Анализ и обработка пропущенных значений (Missing Values)}
    
    \begin{itemize}
        \item Определение доли пропусков в каждом столбце
        \item Анализ паттернов пропусков: являются ли они случайными или системными?
        \item Принятие решения: удаление строк/столбцов, импутация (средним, медианой, моделью), добавление
        флага «был пропуск»
    \end{itemize}
    
    \item \textbf{Анализ распределений (Distribution Analysis)}
    
    \begin{itemize}
        \item Для \textbf{категориальных} признаков построение (\textit{bar plot}) для оценки частоты категорий
        \item Для \textbf{числовых} признаков построение (\textit{histogram}) и (\textit{box plot}) для оценки
        формы распределения, центральной тенденции, разброса и выбросов.
        \item Проверка на нормальность (визуальная или статистические тесты).
    \end{itemize}
    
    \item \textbf{Выявление и анализ выбросов (Outlier Analysis)}
    
    \begin{itemize}
        \item Использование \texttt{box plot}, \texttt{scatter plot} или статистических методов (например,
        правило трёх сигм, интерквартильный размах — IQR).
        \item Принятие решения: удаление, трансформация, каппинг или оставление выбросов (если они несут ценную
        информацию)
    \end{itemize}
    
    \item \textbf{Анализ взаимосвязей (Relationship Analysis)}
    \begin{itemize}
        \item \textbf{Числовые vs Числовые:} Матрица корреляций (Пирсон, Спирмен), диаграмма рассеяния для пар признаков
        \item \textbf{Категориальные vs Числовые:} \textit{Box plot} или \textit{violin plot} для сравнения
        распределений числовой переменной across different categories
        \item \textbf{Категориальные vs Категориальные:} Построение тепловой карты (\textit{heatmap}) сводной
        таблицы или использование критерия хи-квадрат
    \end{itemize}
    
    \item \textbf{Инжиниринг и трансформация признаков (на основе EDA)}
    
    \begin{itemize}
        \item Кодирование категориальных переменных (One-Hot, Label Encoding).
        \item Масштабирование/нормализация числовых признаков (MinMaxScaler, StandardScaler).
        \item Создание новых признаков (например, «площадь на комнату» из «общей площади» и «числа комнат»).
        \item Применение преобразований для нормализации распределения (логарифмирование, корень квадратный).
    \end{itemize}
    
    \item \textbf{Формулировка выводов и гипотез для моделирования}
    \begin{itemize}
        \item Краткое резюме: основные закономерности, обнаруженные проблемы и принятые решения по обработке данных.
        \item Предварительные гипотезы о том, какие признаки могут быть наиболее важными для прогноза target.
    \end{itemize}

\end{enumerate}